\documentclass[10pt]{extarticle}
\usepackage[a4paper,margin=1.5cm]{geometry}
\usepackage[T1]{fontenc}
\usepackage{amsmath,amssymb}
\usepackage{graphicx}
\usepackage[export]{adjustbox}
\usepackage{multicol}
\usepackage{enumitem}
\usepackage{booktabs,array,tabularx}
\setlength{\columnsep}{0.7cm}
\setlength{\parindent}{0pt}
\setlength{\parskip}{3pt}
\pagestyle{empty}
\makeatletter
% figure registry (figdefs.tex is generated by the extract step)
\newcommand{\qfigdef}[4]{\@namedef{qf@file@#1}{#2}\@namedef{qf@w@#1}{#3}\@namedef{qf@wide@#1}{#4}}
\newcommand{\qf@check}[1]{\@ifundefined{qf@file@#1}{\PackageError{score2}{Unknown figure id #1}{}}{}}
\newcommand{\qf@img}[4][]{\includegraphics[width=#2\dimexpr\@nameuse{qf@w@#4}pt\relax,max width=#3,max height=0.45\textheight,#1]{\@nameuse{qf@file@#4}}}
% Full-width questions: finalize wraps every question that contains a wide figure in \qfullbegin...\qfullend,
% so the WHOLE question (stem, figures, options) is set across the page in reading order, then the two columns
% resume. The question's head - its number, stem and first figure (\qfullsplit marks the end of the head) -
% is kept on one page, so a number is never left at a page foot with its figure overleaf; the rest may break
% across pages normally, so no big gaps are left. Without a \qfullsplit, the whole question is kept together
% if it is under 45% of a page. Never write these yourself.
\newif\ifqf@full
\newif\ifqf@boxed
\newsavebox{\qf@box}
\newcommand{\qf@open}{\global\qf@fulltrue\par\global\qf@boxedtrue
  \setbox\qf@box\vbox\bgroup\hsize\textwidth\linewidth\textwidth}
\newcommand{\qf@emit}[1]{\par\egroup\global\qf@boxedfalse
  \ifdim\dimexpr\ht\qf@box+\dp\qf@box\relax<#1\textheight\noindent\usebox{\qf@box}\par\else\unvbox\qf@box\fi}
\newcommand{\qfullsplit}{\ifqf@boxed\qf@emit{0.85}\fi}
\newcommand{\qf@close}{\ifqf@boxed\qf@emit{0.45}\fi\global\qf@fullfalse}
\newcommand{\qfullbegin}{\end{multicols}\qf@open}
\newcommand{\qfullend}{\qf@close\begin{multicols}{2}}
\newcommand{\qfullstart}{\qf@open}
\newcommand{\qfulldone}{\qf@close}
% \qfig[scale]{id}: block figure, centred, directly where it is written (size set by the tool).
% Fallback if a wide figure is not inside a full-width question: break out of the columns in place - no floats.
\newcommand{\qfig}[2][1]{\qf@check{#2}%
  \ifqf@full\qf@blk{#1}{#2}\else
  \ifnum\@nameuse{qf@wide@#2}=1 \expandafter\@firstoftwo\else\expandafter\@secondoftwo\fi
  {\end{multicols}\begin{center}\qf@img{#1}{\textwidth}{#2}\end{center}\begin{multicols}{2}}%
  {\qf@blk{#1}{#2}}\fi}
\newcommand{\qf@blk}[2]{\par\smallskip{\centering\qf@img{#1}{\linewidth}{#2}\par}\smallskip}
% \qfigi[scale]{id}: inline figure (inside option lists / tables / minipages); never breaks columns.
\newcommand{\qfigi}[2][1]{\qf@check{#2}\qf@img[valign=c]{#1}{\linewidth}{#2}}
\makeatother
% structure macros - lint reads these to check completeness and figure placement
\newcommand{\papersec}[1]{\par\bigskip{\centering\Large\bfseries PAPER-#1\par}\medskip}
\newcommand{\partsec}[2]{\par\medskip\noindent{\large\bfseries PART-#1: #2}\par\smallskip}
\newcommand{\sect}[1]{\par\smallskip\noindent\textbf{#1}\par}
\newcommand{\qn}[1]{\par\smallskip\noindent\textbf{#1.}~}
\newcommand{\paraq}[2]{\par\smallskip\noindent\textbf{Paragraph for Questions #1 and #2}\par}
\newenvironment{qoptions}{\begin{itemize}[leftmargin=1.8em,itemsep=1pt,topsep=2pt,label={}]}{\end{itemize}}
\newcommand{\opt}[2]{\item[(#1)] #2}
% figure registry (generated by score2_tools.py finalize)
\qfigdef{p1_phy_s1_q01_1}{Score Advanced - Test 06 (04 Sep 2021)-img/p1_phy_s1_q01_1.png}{102.7}{0}
\qfigdef{p1_phy_s1_q02_1}{Score Advanced - Test 06 (04 Sep 2021)-img/p1_phy_s1_q02_1.png}{116.3}{0}
\qfigdef{p1_phy_s1_q03_1}{Score Advanced - Test 06 (04 Sep 2021)-img/p1_phy_s1_q03_1.png}{78.6}{0}
\qfigdef{p1_phy_s1_q04_1}{Score Advanced - Test 06 (04 Sep 2021)-img/p1_phy_s1_q04_1.png}{128.1}{0}
\qfigdef{p1_phy_s2_para9_10_1}{Score Advanced - Test 06 (04 Sep 2021)-img/p1_phy_s2_para9_10_1.png}{136.8}{0}
\qfigdef{p1_phy_s3_q01_1}{Score Advanced - Test 06 (04 Sep 2021)-img/p1_phy_s3_q01_1.png}{107.0}{0}
\qfigdef{p1_phy_s3_q04_1}{Score Advanced - Test 06 (04 Sep 2021)-img/p1_phy_s3_q04_1.png}{76.4}{0}
\qfigdef{p1_phy_s3_q05_1}{Score Advanced - Test 06 (04 Sep 2021)-img/p1_phy_s3_q05_1.png}{118.8}{0}
\qfigdef{p1_phy_s3_q06_1}{Score Advanced - Test 06 (04 Sep 2021)-img/p1_phy_s3_q06_1.png}{47.2}{0}
\qfigdef{p1_phy_s3_q07_1}{Score Advanced - Test 06 (04 Sep 2021)-img/p1_phy_s3_q07_1.png}{122.0}{0}
\qfigdef{p1_phy_s3_q08_1}{Score Advanced - Test 06 (04 Sep 2021)-img/p1_phy_s3_q08_1.png}{135.2}{0}
\qfigdef{p1_chem_s1_q01_1}{Score Advanced - Test 06 (04 Sep 2021)-img/p1_chem_s1_q01_1.png}{316.2}{1}
\qfigdef{p1_chem_s1_q02_A}{Score Advanced - Test 06 (04 Sep 2021)-img/p1_chem_s1_q02_A.png}{44.9}{0}
\qfigdef{p1_chem_s1_q02_B}{Score Advanced - Test 06 (04 Sep 2021)-img/p1_chem_s1_q02_B.png}{46.8}{0}
\qfigdef{p1_chem_s1_q02_C}{Score Advanced - Test 06 (04 Sep 2021)-img/p1_chem_s1_q02_C.png}{46.1}{0}
\qfigdef{p1_chem_s1_q02_D}{Score Advanced - Test 06 (04 Sep 2021)-img/p1_chem_s1_q02_D.png}{47.0}{0}
\qfigdef{p1_chem_s1_q04_1}{Score Advanced - Test 06 (04 Sep 2021)-img/p1_chem_s1_q04_1.png}{208.2}{0}
\qfigdef{p1_chem_s2_para7_8_1}{Score Advanced - Test 06 (04 Sep 2021)-img/p1_chem_s2_para7_8_1.png}{235.4}{0}
\qfigdef{p1_chem_s2_q07_A}{Score Advanced - Test 06 (04 Sep 2021)-img/p1_chem_s2_q07_A.png}{68.9}{0}
\qfigdef{p1_chem_s2_q07_B}{Score Advanced - Test 06 (04 Sep 2021)-img/p1_chem_s2_q07_B.png}{47.3}{0}
\qfigdef{p1_chem_s2_q07_C}{Score Advanced - Test 06 (04 Sep 2021)-img/p1_chem_s2_q07_C.png}{72.3}{0}
\qfigdef{p1_chem_s2_q07_D}{Score Advanced - Test 06 (04 Sep 2021)-img/p1_chem_s2_q07_D.png}{59.2}{0}
\qfigdef{p1_chem_s2_q08_1}{Score Advanced - Test 06 (04 Sep 2021)-img/p1_chem_s2_q08_1.png}{169.5}{0}
\qfigdef{p1_chem_s2_q08_A}{Score Advanced - Test 06 (04 Sep 2021)-img/p1_chem_s2_q08_A.png}{39.0}{0}
\qfigdef{p1_chem_s2_q08_B}{Score Advanced - Test 06 (04 Sep 2021)-img/p1_chem_s2_q08_B.png}{25.8}{0}
\qfigdef{p1_chem_s2_q08_C}{Score Advanced - Test 06 (04 Sep 2021)-img/p1_chem_s2_q08_C.png}{42.6}{0}
\qfigdef{p1_chem_s2_q08_D}{Score Advanced - Test 06 (04 Sep 2021)-img/p1_chem_s2_q08_D.png}{40.4}{0}
\qfigdef{p1_chem_s3_q06_1}{Score Advanced - Test 06 (04 Sep 2021)-img/p1_chem_s3_q06_1.png}{121.8}{0}
\qfigdef{p2_phy_s1_q01_1}{Score Advanced - Test 06 (04 Sep 2021)-img/p2_phy_s1_q01_1.png}{63.9}{0}
\qfigdef{p2_phy_s1_q02_1}{Score Advanced - Test 06 (04 Sep 2021)-img/p2_phy_s1_q02_1.png}{125.1}{0}
\qfigdef{p2_phy_s1_q05_1}{Score Advanced - Test 06 (04 Sep 2021)-img/p2_phy_s1_q05_1.png}{119.0}{0}
\qfigdef{p2_phy_s1_q06_1}{Score Advanced - Test 06 (04 Sep 2021)-img/p2_phy_s1_q06_1.png}{100.2}{0}
\qfigdef{p2_phy_s2_q07_1}{Score Advanced - Test 06 (04 Sep 2021)-img/p2_phy_s2_q07_1.png}{91.4}{0}
\qfigdef{p2_phy_s2_q08_1}{Score Advanced - Test 06 (04 Sep 2021)-img/p2_phy_s2_q08_1.png}{101.0}{0}
\qfigdef{p2_phy_s2_q09_1}{Score Advanced - Test 06 (04 Sep 2021)-img/p2_phy_s2_q09_1.png}{158.5}{0}
\qfigdef{p2_phy_s2_q10_1}{Score Advanced - Test 06 (04 Sep 2021)-img/p2_phy_s2_q10_1.png}{78.7}{0}
\qfigdef{p2_phy_s2_q10_2}{Score Advanced - Test 06 (04 Sep 2021)-img/p2_phy_s2_q10_2.png}{236.8}{0}
\qfigdef{p2_phy_s3_q02_1}{Score Advanced - Test 06 (04 Sep 2021)-img/p2_phy_s3_q02_1.png}{154.5}{0}
\qfigdef{p2_phy_s3_q03_1}{Score Advanced - Test 06 (04 Sep 2021)-img/p2_phy_s3_q03_1.png}{179.9}{0}
\qfigdef{p2_phy_s3_q04_1}{Score Advanced - Test 06 (04 Sep 2021)-img/p2_phy_s3_q04_1.png}{87.5}{0}
\qfigdef{p2_phy_s3_q07_1}{Score Advanced - Test 06 (04 Sep 2021)-img/p2_phy_s3_q07_1.png}{156.0}{0}
\qfigdef{p2_chem_s1_q04_A}{Score Advanced - Test 06 (04 Sep 2021)-img/p2_chem_s1_q04_A.png}{84.9}{0}
\qfigdef{p2_chem_s1_q04_B}{Score Advanced - Test 06 (04 Sep 2021)-img/p2_chem_s1_q04_B.png}{133.9}{0}
\qfigdef{p2_chem_s1_q04_C}{Score Advanced - Test 06 (04 Sep 2021)-img/p2_chem_s1_q04_C.png}{148.0}{0}
\qfigdef{p2_chem_s1_q04_D}{Score Advanced - Test 06 (04 Sep 2021)-img/p2_chem_s1_q04_D.png}{128.5}{0}
\qfigdef{p2_chem_s2_q07_1}{Score Advanced - Test 06 (04 Sep 2021)-img/p2_chem_s2_q07_1.png}{54.9}{0}
\qfigdef{p2_chem_s2_q07_2}{Score Advanced - Test 06 (04 Sep 2021)-img/p2_chem_s2_q07_2.png}{36.8}{0}
\qfigdef{p2_chem_s2_q07_3}{Score Advanced - Test 06 (04 Sep 2021)-img/p2_chem_s2_q07_3.png}{66.3}{0}
\qfigdef{p2_chem_s2_q07_4}{Score Advanced - Test 06 (04 Sep 2021)-img/p2_chem_s2_q07_4.png}{79.1}{0}
\qfigdef{p2_chem_s2_q09_1}{Score Advanced - Test 06 (04 Sep 2021)-img/p2_chem_s2_q09_1.png}{111.5}{0}
\qfigdef{p2_chem_s2_q10_1}{Score Advanced - Test 06 (04 Sep 2021)-img/p2_chem_s2_q10_1.png}{126.8}{0}
\qfigdef{p2_chem_s2_q10_2}{Score Advanced - Test 06 (04 Sep 2021)-img/p2_chem_s2_q10_2.png}{148.2}{0}
\qfigdef{p2_chem_s2_q10_3}{Score Advanced - Test 06 (04 Sep 2021)-img/p2_chem_s2_q10_3.png}{146.2}{0}
\qfigdef{p2_chem_s2_q10_4}{Score Advanced - Test 06 (04 Sep 2021)-img/p2_chem_s2_q10_4.png}{109.4}{0}
\qfigdef{p2_chem_s3_q07_1}{Score Advanced - Test 06 (04 Sep 2021)-img/p2_chem_s3_q07_1.png}{245.2}{0}
\raggedcolumns
\makeatletter\@beginparpenalty=10000\makeatother
\newcolumntype{L}[1]{>{\raggedright\arraybackslash}p{#1}}
\begin{document}
\begin{center}{\Large\bfseries Score Advanced --- Test 06 (04 Sep 2021)}\end{center}
\begin{multicols}{2}
\raggedcolumns
\papersec{1}
\partsec{1}{PHYSICS}
\sect{SECTION-I (i)}
\qn{1} A spherical surface separates air \& medium for which $\mu$ = 1.615 for violet and $\mu$ = 1.600 for red color. A paraxial beam parallel to optic axis is incident on the surface as shown. The distance between point of convergence for violet and red color is $\Delta$f.\par
\qfig{p1_phy_s1_q01_1}
\begin{qoptions}
\opt{A}{$\Delta$f = 0.41 cm}
\opt{B}{Point of convergence for red is closer to optical centre than that for violet.}
\opt{C}{Point of convergence for violet is closer to optical centre than that for red.}
\opt{D}{$\Delta$f = 0.84 cm.}
\end{qoptions}
\qn{2} A cubic frame is made of 12 rods each of mass m and length $\ell$.\par
\qfig{p1_phy_s1_q02_1}
\begin{qoptions}
\opt{A}{The moment of inertia of the cube about its face diagonal (BD) is $\dfrac{23}{3}$m$\ell^2$.}
\opt{B}{The moment of inertia of the cube about an axis passing through the center of cube and perpendicular to one of the face (HI) is $\dfrac{14}{3}$m$\ell^2$.}
\opt{C}{The moment of inertia of the cube about one of the side of cube (AB) is $\dfrac{32}{3}$m$\ell^2$.}
\opt{D}{The moment of inertia of the cube about an axis passing through the center of the face and the midpoint of the one of the side (JK) is $\dfrac{23}{6}$m$\ell^2$.}
\end{qoptions}
\qn{3} A smooth sphere of mass 1 kg is tied to a fixed point by an inextensible string and suspended from the roof. Another identical sphere impinges directly on it with speed 5 m/sec in a direction making an acute angle 30$^\circ$ with the string. Co-efficient of restitution is $\dfrac{1}{2}$.\par
\qfig{p1_phy_s1_q03_1}
\begin{qoptions}
\opt{A}{Speed of sphere B just after collision is 3 m/s.}
\opt{B}{Speed of sphere A just after collision is 2 m/s.}
\opt{C}{Impulse due to normal force between two spheres during collision is 6 N-S}
\opt{D}{Impulse due to tension in string during collision is 3 N-S.}
\end{qoptions}
\qn{4} Rank the magnitude of $\oint\vec{B}\cdot d\vec{\ell}$ for the closed paths shown in figure from the smallest to largest.\par
\qfig{p1_phy_s1_q04_1}
\begin{qoptions}
\opt{A}{a, b, c, d}
\opt{B}{a, b, d, c}
\opt{C}{a, d, c, b}
\opt{D}{a, c, b, d}
\end{qoptions}
\qn{5} A point charge q is placed at origin. The electrostatic energy stored in the space around the point charge in various volumes as mentioned in listed in options below. Choose the correct option(s) below.
\begin{qoptions}
\opt{A}{A small cube with edge size a (a $<<$ r) placed at (r, 0, 0) has $\dfrac{q^2a^3}{32\pi^2\varepsilon_0 r^4}$ amount of energy.}
\opt{B}{A small cube with edge size a (a $<<$ r) placed at (r, 0, 0) has $\dfrac{q^2a^3}{16\pi^2\varepsilon_0 r^4}$ amount of energy}
\opt{C}{Energy stored in volume given by a$^2$ $\leq$ x$^2$ + y$^2$ + z$^2$ $\leq$ 4a$^2$ is $\dfrac{q^2}{16\pi\varepsilon_0 a}$}
\opt{D}{Energy stored in volume given by a$^2$ $\leq$ x$^2$ + y$^2$ + z$^2$ $\leq$ 4a$^2$ is $\dfrac{q^2}{32\pi\varepsilon_0 a}$}
\end{qoptions}
\qn{6} A string of mass 0.2 kg and length 2m is tied at two ends to fixed supports under a tension of 10 N. A point P on the string is found to travel from one extreme to other in 0.1s. Taking one end as x = 0 and the other end x = 2 m and t = 0 as the time when P is at rest. (Position of P is x) The correct statements will be :-
\begin{qoptions}
\opt{A}{For 0 $<$ t $<$ 0.1s, energy flows across P in positive x-direction for 0 $<$ x $<$ 1 m}
\opt{B}{For 0 $<$ t $<$ 0.05 s, energy flows across P in negative x-direction for 0 $<$ x $<$ 1m}
\opt{C}{At t = 0.05 s, rate of energy flow through P is zero for x = 0.5 m}
\opt{D}{At t = 0.1s, rate of energy flow through P for all values of x is zero.}
\end{qoptions}
\sect{SECTION-I (ii)}
\paraq{7}{8}
Nuclear binding energy varies from nucleus to nucleus, it is necessary to compare the binding energy per nucleon basis. If we find binding energy per nucleon for a given element it can be given as binding energy divided by the nucleon number A as\par
$(\text{BE})_N = \dfrac{\Delta E}{A} = \dfrac{\Delta mc^2}{A}$\par
We can define binding energy per nucleon theoretically as the amount of energy needed to remove a nucleon from nucleus of an element.
\qn{7} The binding energy of deuteron $\left({}^{2}_{1}\text{H}\right)$ is 1.15 MeV per nucleon and an alpha particle $\left({}^{4}_{2}\text{He}\right)$ has a binding energy of 7.1 MeV per nucleon. Then in the reaction ${}^{2}_{1}\text{H}$ + ${}^{2}_{1}\text{H}$ $\rightarrow$ ${}^{4}_{2}\text{He}$ + Q the energy Q released is :
\begin{qoptions}
\opt{A}{1MeV}
\opt{B}{11.9 MeV}
\opt{C}{23.8 MeV}
\opt{D}{931 MeV}
\end{qoptions}
\qn{8} What is the energy released in the fission reaction ${}^{236}_{92}$U $\rightarrow$ ${}^{117}_{46}$X + ${}^{117}_{46}$Y + 2 ${}^{1}_{0}$n given that the binding energy per nucleon of X and Y is 8.5 MeV and that of ${}^{236}_{92}$U is 7.6 MeV?
\begin{qoptions}
\opt{A}{176.2 MeV}
\opt{B}{183.7 MeV}
\opt{C}{195.4 MeV}
\opt{D}{208.6 MeV}
\end{qoptions}
\paraq{9}{10}
An electromagnetic wave can be represented by E = A sin (kx -- $\omega$t + $\phi$), where E is electric field associated with wave. According this equation, for any value of x, E remains sinusoidal for $-\infty$ $<$ t $<$ $\infty$. Obviously this corresponds to an idealised situation because radiation from ordinary source consists of finite size wavetrains. In general, electric field remains sinusoidal only for times of order $\tau_C$, which is called coherence time. In simpler language it means that for times of order $\tau_C$, a wave will have a definite phase. The finite value of coherence time could be due to many factors, for example if radiating atom undergoes collision with another atom then wave train undergoes an abrupt phase change or due to the fact that an atom responsible for emitting radiation has a finite life time in the energy level from which it drops to lower energy level, while radiating. Concept of coherence time can be easily understood using young's double slit experiment. Let interference pattern is observed around point P at time t, due to superposition of waves emanting from S$_1$ and S$_2$ at times t = $\dfrac{r_1}{c}$ and $\dfrac{r_2}{c}$ respectively, where r$_1$ and r$_2$ are the distance S$_1$P \& S$_2$P. Obviously if $\dfrac{r_2-r_1}{c}$ $<<$ $\tau_c$, \{where c = 3 $\times$ 10$^8$ m/s\} then, wavetrain arriving at point P from S$_1$ \& S$_2$ will have a definite phase relationship and an interference pattern of good contrast will be obtained.\par
\qfig{p1_phy_s2_para9_10_1}
\qn{9} Which of the following statements is incorrect in the diagram given :-
\begin{qoptions}
\opt{A}{Central fringe (near O) will always have a good contrast.}
\opt{B}{As we move on the screen towards higher orders, contrast will become poorer.}
\opt{C}{If we move S, parallel to S$_1$S$_2$ there will be no change in contrast of fringes on screen around O.}
\opt{D}{If we move on screen away from O, fringes will disappear beyond a certain limit.}
\end{qoptions}
\qn{10} If coherence time is of order 10$^{-10}$ second and screen is placed at a very large distance from slits in the given figure, then :-
\begin{qoptions}
\opt{A}{For d = 0.3 cm, we will not get any fringes on screen (except central maxima)}
\opt{B}{For d = 0.03 cm we will get fringes on a large part of screen.}
\opt{C}{For d = 3 cm we will get fringes on entire screen.}
\opt{D}{We will get fringes on entire screen, irrespective of value of d.}
\end{qoptions}
\sect{SECTION-II}
\qn{1} A rod of length 1m rotates about one of its end points with an angular velocity 2 rad/sec in a plane perpendicular to the magnetic field B = 2T as shown in the figure. Then find magnitude of electric field (in SI unit) at the mid point of the rod.\par
\qfig{p1_phy_s3_q01_1}
\qn{2} A steel rod at 25$^\circ$C is bolted at both ends and then cooled. By how many $^\circ$C should the rod be cooled so that it will rupture ? Assume that till rupture, Hooke's law is obeyed. ($\alpha_{\text{steel}}$ = 10 $\times$ 10$^{-6}$ $^\circ$C, Y = 2 $\times$ 10$^{11}$ N/m$^2$ and $\sigma_b$ (breaking stress of steel rod) = 4 $\times$ 10$^8$ N/m$^2$)
\qn{3} The velocity distribution in a viscous flow over a plate is given by u = 4y -- y$^2$ for y $\leq$ 2m where u = velocity in m/s at a point distant y from the plate. If the coefficient of dynamic viscosity is 1.5 Pa.s. Determine the shear stress at y = 1m (SI units)
\qn{4} In the figure shown water is filled in the cylindrical tank of cross-sectional area 100 cm$^2$ up to a height of 10 cm. The mass of the thin, tight fitting and frictionless piston is 100 kg. The small hole in the piston has area $2\sqrt{2}$ mm$^2$. Due to weight of the piston water spurts upwards in a jet from the hole. Find the time (in sec.) in which the tank becomes empty. (Assume that spurted water does not fall back on piston)\par
\qfig{p1_phy_s3_q04_1}
\qn{5} A battery of emf E$_0$ = 12.5 V is connected across a 4m long uniform wire having resistance 4$\Omega$/m. Two cell of small EMFs $\xi_1$ = 2V and $\xi_2$ = 4V having internal resistance 2$\Omega$ and 6$\Omega$ respectively are connected as shown in the figure. If galvanometer shows no deflection at the point N, the distance of point N (in m) from the point A is equal to :\par
\qfig{p1_phy_s3_q05_1}
\qn{6} If the block shown in figure is displaced slightly and released, find the time period of oscillations of the block is found to be $2\pi\sqrt{\dfrac{\text{Pm}}{\text{K}}}$. Find P.\par
\qfig{p1_phy_s3_q06_1}
\qn{7} A boat travel, between two points A and B on the opposite banks of a river (figure), always following the line AB. The distance S between points A and B is 240 m. The velocity of the river flow v$_r$ = $2\sqrt{2}$ m/s is constant over the entire width of the river. The line AB makes an $\alpha$ = 45$^\circ$ with the direction of the current. The boat moves to cover the distance AB and back in a time t = $2\sqrt{3}$ min. The angle $\beta$ remains the same during the passage from A to B and from B to A. With what velocity u$_b$ should the boat travel ?\par
\qfig{p1_phy_s3_q07_1}
\qn{8} At an instant t = 0, Block A of mass 1 kg is moving with speed 8 m/s towards right on rough surface of block B of mass 3 kg. Block B, which is placed on smooth horizontal surface is moving with speed 4 m/s towards right at same instant (t = 0). The net work done against the frictional force in long time is (block A will not fall from block B)\par
\qfig{p1_phy_s3_q08_1}
\qfullbegin
\partsec{2}{CHEMISTRY}
\sect{SECTION-I (i)}
\qn{1} \par
\qfig{p1_chem_s1_q01_1}
\qfullsplit
product. (X + Y) is :
\begin{qoptions}
\opt{A}{3}
\opt{B}{2}
\opt{C}{4}
\opt{D}{8}
\end{qoptions}
\qfullend
\qn{2} Which is best graphical representation of atomic radius for 3d series elements.
\begin{qoptions}
\opt{A}{\qfigi{p1_chem_s1_q02_A}}
\opt{B}{\qfigi{p1_chem_s1_q02_B}}
\opt{C}{\qfigi{p1_chem_s1_q02_C}}
\opt{D}{\qfigi{p1_chem_s1_q02_D}}
\end{qoptions}
\qn{3} Which is/are correct statements ?
\begin{qoptions}
\opt{A}{AgBr shows Frenkel defect as well as schottky defect.}
\opt{B}{Density may be increase due to impurity defect.}
\opt{C}{End centred cubic unit cell does not exist in nature}
\opt{D}{Orthorhombic crystal system has three bravias lattices.}
\end{qoptions}
\qn{4} \par
\qfig{p1_chem_s1_q04_1}
Which of the reactant will formed optically active product after reductive ozonolysis.
\begin{qoptions}
\opt{A}{I}
\opt{B}{II}
\opt{C}{III}
\opt{D}{IV}
\end{qoptions}
\qn{5} Which process of metal purification is also associated with distillation of metal
\begin{qoptions}
\opt{A}{Parke's process}
\opt{B}{Amalgamation}
\opt{C}{Polling}
\opt{D}{Puddling process}
\end{qoptions}
\qn{6} Acetone and carbon disulphide form binary liquid solution showing positive deviation from Raoult's law with a maxima in P$_T$ vs X graph. The normal boiling point (T$_b$) of pure acetone is less than that of pure CS$_2$. Pick out the incorrect statements among the following.
\begin{qoptions}
\opt{A}{Boiling temperature of mixture is always less than boiling temperature of acetone.}
\opt{B}{Boiling temperature of Azeotropic mixture is less than boiling temperature of pure CS$_2$}
\opt{C}{When a small amount CS$_2$ is added to excess of acetone boiling point of Resulting mixture increases.}
\opt{D}{A mixture of CS$_2$ and CH$_3$COCH$_3$ can be completely separated by simple fractional distillation.}
\end{qoptions}
\sect{SECTION-I (ii)}
\paraq{7}{8}
Observe the following reaction sequence and give answers to questions given below :\par
\qfig{p1_chem_s2_para7_8_1}
\qn{7} Identify (Y) :
\begin{qoptions}
\opt{A}{\qfigi{p1_chem_s2_q07_A}}
\opt{B}{\qfigi{p1_chem_s2_q07_B}}
\opt{C}{\qfigi{p1_chem_s2_q07_C}}
\opt{D}{\qfigi{p1_chem_s2_q07_D}}
\end{qoptions}
\qn{8} \qfig{p1_chem_s2_q08_1}
\begin{qoptions}
\opt{A}{\qfigi{p1_chem_s2_q08_A}}
\opt{B}{\qfigi{p1_chem_s2_q08_B}}
\opt{C}{\qfigi{p1_chem_s2_q08_C}}
\opt{D}{\qfigi{p1_chem_s2_q08_D}}
\end{qoptions}
\paraq{9}{10}
The atomic orbitals combine to form new set of equivalent orbitals known as hybrid orbitals. Unlike pure orbitals, the hybrid orbitals are used in bond formation. The phenomenon is known as hybridisation which can be defined as the process of intermixing of the orbitals of slightly different energies in the formation of new set of orbitals of equivalent energies.
\qn{9} Which molecule is non planar as well as polar
\begin{qoptions}
\opt{A}{IF}
\opt{B}{IF$_3$}
\opt{C}{IF$_5$}
\opt{D}{IF$_7$}
\end{qoptions}
\qn{10} Which molecule is having same electronic geometry as well as molecular geometry
\begin{qoptions}
\opt{A}{NCl$_3$}
\opt{B}{PCl$_5$}
\opt{C}{SF$_4$}
\opt{D}{XeF$_4$}
\end{qoptions}
\sect{SECTION-II}
\qn{1} Pt, H$_2$(1 bar) $|$ 1 M NaA(aq) $\|$ 10$^{-x}$ M HCl (aq) $|$ H$_2$(2 bar), Pt\par
If emf of above cell is 0.441. What is the value of x ?\par
(Given : 2.303 RT/F = 0.06 ; log2 = 0.3) (pK$_a$ (HA) = 5)
\qn{2} How many of the following statements regarding kinetics of a chemical reaction are correct.\par
(i) The probability of collision of three or more molecules with proper orientation at same time is very rare thats why reactions having molecularity three or more are very rare.\par
(ii) In a second order reaction the number of reactant can be one but the rate will be proportional to the square of its concentration.\par
(iii) The rate of a chemical reaction at any time is always proportional to the initial concentration of the reactant molecules raised to the power equal to their order.\par
(iv) The rate of a chemical reaction also depends upon the strength of bond of reactant molecules and the number of bond to be broken.\par
(v) The total kinetic energy of the colliding molecules must be greater than some minimum value.\par
(vi) The colliding particles must be properly oriented in space when they collide.
\qn{3} x = number of optically pure aldehyde or ketone having molecular formula of C$_5$H$_{10}$O give chiral alcohol but not racemic mixture with NaBH$_4$ / CH$_3$--OH. (Consider enantiomer separately).
\qn{4} How many of the following compound(s) give precipitate with solution of chrome alum.\par
(i) BaCl$_2$ \quad (ii) Na$_2$CO$_3$ \quad (iii) (NH$_4$)$_2$S \quad (iv) Excess NaOH
\qn{5} 1 mole of an ideal gas initially present in a 2L insulated cylinder at 300 K is allowed to expand against vaccum to 8 L. then $|\Delta$G$_{\text{sys}}|$ in calories will be [ln2 = 0.7]
\qn{6} \qfig{p1_chem_s3_q06_1}
Number of CH$_3$ -- I consumed during the following reaction. In which Hoffmann exhaustive elimination take place.
\qn{7} For balanced chemical reaction : `X'NH$_3$(g)$_{\text{excess}}$ + `Y'Cl$_2$O $\xrightarrow{\text{Warm}}$ ...........+ 6 NH$_4$Cl + .............\par
find the value of expression $|$X -- Y$|$.\par
Where, `X' and `Y' are number of moles of NH$_3$ and Cl$_2$O respectively required to produce six moles of NH$_4$Cl.
\qn{8} One mole of a real gas obeys P(Vm -- b) = RT ; where `b' and `R' are constants with usual meanings. If occupied volume of the gas is `10 b' then the compressibility factor of gas will be :
\partsec{3}{MATHEMATICS}
\sect{SECTION-I (i)}
\qn{1} Consider the planes P$_1$ : 2x + y + z + 4 = 0, P$_2$ : y -- z + 4 = 0 and P$_3$ : 3x + 2y + z + 8 = 0. Let L$_1$, L$_2$, L$_3$ be the lines of intersection of the planes P$_2$ and P$_3$, P$_3$ and P$_1$, and P$_1$ and P$_2$ respectively. Then,
\begin{qoptions}
\opt{A}{Atleast two of the lines L$_1$, L$_2$ and L$_3$ are non-parallel}
\opt{B}{Atleast two of the lines L$_1$, L$_2$ and L$_3$ are parallel}
\opt{C}{The three planes intersect in a line}
\opt{D}{The three planes form a triangular prism}
\end{qoptions}
\qn{2} Let f(x) be a differentiable function for x $>$ --1 and f(0) = 1, f$'$(0) $>$ 0 also f$'$(x)(1 + f(x)) = 1 + x then
\begin{qoptions}
\opt{A}{f$'$(x) $>$ 0 $\forall$ x $>$ 0}
\opt{B}{f$'$(x) $<$ 0 $\forall$ x $>$ 0}
\opt{C}{f(x) $>$ --1 $\forall$ x $>$ 0}
\opt{D}{f(x) $<$ --1 $\forall$ x $>$ 0}
\end{qoptions}
\qn{3} If the quadratic equation x$^2$ + 2ax + b(a -- 1) = 0 has real roots for all real values of a and b, then which of the following is/are correct ?
\begin{qoptions}
\opt{A}{Greatest value of b is 4}
\opt{B}{No such value of b is possible}
\opt{C}{Least value of b is 0}
\opt{D}{b can be any integer}
\end{qoptions}
\qn{4} Suppose three real numbers a, b, c are in G.P. Let z = $\dfrac{a+ib}{c-ib}$, then
\begin{qoptions}
\opt{A}{z = $\dfrac{ib}{c}$}
\opt{B}{z = $\dfrac{ia}{b}$}
\opt{C}{z = $\dfrac{ia}{c}$}
\opt{D}{z = 0}
\end{qoptions}
\qn{5} In $\Delta$ABC, if $\angle$B = $\sec^{-1}\left(\dfrac{5}{4}\right)$ + $\text{cosec}^{-1}\sqrt{5}$, $\angle$C = $\text{cosec}^{-1}\left(\dfrac{25}{7}\right)$ + $\cot^{-1}\left(\dfrac{9}{13}\right)$ and c = 3. Which of the following statement(s) is/are correct ? (where all symbols used have their usual meanings in a triangle)
\begin{qoptions}
\opt{A}{tan A, tan B, tan C are in arithmetic progression.}
\opt{B}{Inradius the triangle is irrational}
\opt{C}{Area of $\Delta$ABC is rational.}
\opt{D}{Semiperimeter is rational}
\end{qoptions}
\qn{6} Let $\displaystyle\int_0^{\pi/2^n}\ln(\sin x)dx = a_n$, $\displaystyle\int_0^{\pi/2^n}\ln(\cos x)dx = b_n$, then
\begin{qoptions}
\opt{A}{a$_1$ + b$_1$ = $-\pi$ ln 2}
\opt{B}{a$_2$ + b$_2$ = $-\dfrac{\pi}{2}$ ln 2}
\opt{C}{a$_2$ + b$_2$ = $\dfrac{\pi}{2}$ ln 2}
\opt{D}{a$_1$ + b$_1$ = $-\dfrac{\pi}{4}$ ln 2}
\end{qoptions}
\sect{SECTION-I (ii)}
\paraq{7}{8}
Let f(x) be a polynomial with positive leading coefficient satisfying f(0) = 0 \& $f(f(x)) = x\displaystyle\int_0^x f(t)\,dt$ $\forall$ x $\in$ R.
\qn{7} Two perpendicular tangents to the curve y = f(x) will intersect on the curve/line
\begin{qoptions}
\opt{A}{x$^2$ + y$^2$ = 3}
\opt{B}{y = $\dfrac{1}{4\sqrt{3}}$}
\opt{C}{y = $-\dfrac{\sqrt{3}}{4}$}
\opt{D}{y$^2$ = $\dfrac{3x}{-16}$}
\end{qoptions}
\qn{8} If a line having slope 1 and passing through ($\alpha$, 0) neither cuts nor touches the curve y = f(x), then the value of $\alpha$ must lie in the interval.
\begin{qoptions}
\opt{A}{(0, $\infty$)}
\opt{B}{$\left(\dfrac{\sqrt{3}}{4},\infty\right)$}
\opt{C}{$\left(\dfrac{3-\sqrt{3}}{4},\infty\right)$}
\opt{D}{$\left(-\dfrac{3}{4},\infty\right)$}
\end{qoptions}
\paraq{9}{10}
Let A and B are two square matrix of order 2 $\times$ 2, where det(A) = 1, det(B) = 2, then
\qn{9} Find det(A + $\alpha$B) -- det($\alpha$A+ B), $\alpha\in$ R :
\begin{qoptions}
\opt{A}{$\alpha^2$}
\opt{B}{0}
\opt{C}{$\alpha^2 - 1$}
\opt{D}{None of these}
\end{qoptions}
\qn{10} Find det(A -- $\alpha$B) -- det(--$\alpha$A + B) , $\alpha\in$ R :
\begin{qoptions}
\opt{A}{$\alpha^2$}
\opt{B}{0}
\opt{C}{$\alpha^2 - 1$}
\opt{D}{None of these}
\end{qoptions}
\sect{SECTION-II}
\qn{1} Find the distance of the point $\hat{i}+2\hat{j}+3\hat{k}$ from the plane $\vec{r}\cdot\left(\hat{i}+\hat{j}+\hat{k}\right) = 5$ measured parallel to the vector $2\hat{i}+3\hat{j}-6\hat{k}$.
\qn{2} Number of values of `p' for which the equation, (p$^2$ -- 3p + 2)x$^2$ -- (p$^2$ -- 5p + 4)x + p -- p$^2$ = 0 possess more than two roots, is ?
\qn{3} Area of triangle formed by the asymptotes and any tangent to $\dfrac{x^2}{a^2}-\dfrac{y^2}{b^2}=1$ is 36 sq. units, where a \& b are integers, then number of such hyperbolas is ?
\qn{4} From an external point P, two tangents are drawn to the parabola x$^2$ = --8y. If chord of contact of these tangents is tangential to the circle x$^2$ + y$^2$ = 4, then locus of P will be a hyperbola length of whose transverse axis is ?
\qn{5} All the face cards from a pack of 52 playing cards are removed. From the remaining pack half of the cards are randomly removed whithout looking at them and then randomly drawn two cards simultaneously from the remaining. If the probability that, two cards drawn are both aces, is $\dfrac{p\left({}^{38}C_{20}\right)}{{}^{40}C_{20}\cdot{}^{20}C_2}$, then the value of p is ?
\qn{6} Normal to the curve y = f(x) at (1, 1) is 3x + 4y = 7. Also $f(x)\geqslant\dfrac{4x-1}{3}$ for x $\geq$ 1 and $f(x)\leqslant\dfrac{4x-1}{3}$ for x $\leq$ 1. where f(x) is twice differentiable everywhere in its domain. Evaluate $\displaystyle\lim_{x\to1}\dfrac{3f'(x)-2f(x)-2x}{f(x)-x^2}$.
\qn{7} If a, b, c be three natural numbers in A.P. and a + b + c = 30, then the possible number of values of a, b, c is ?
\qn{8} If the co-efficient of x$^4$ in the expansion of (1 -- x + 2x$^2$)$^{12}$ is ${}^{n+2}C_{r+1}$ + r $\cdot$ ${}^{n+1}C_r$ + ${}^nC_r$, then the coefficient of x$^r$ in the expansion of (1 + x)$^{n/2}$ is
\raggedcolumns
\papersec{2}
\partsec{1}{PHYSICS}
\sect{SECTION-I (i)}
\qn{1} Two identical cylindrical tanks are connected by a narrow tube with a cork at its middle (fig). The radius of a tank is R = 20.0 cm, the radius of the tube is r = 1.0 mm. The length of the tube is $\ell$ = 1m. Water ($\rho$ = 1000 kg/m$^3$) is poured into one of the tanks to a height of h = 50 cm, while the second tank is empty. At the instant t = 0, the cork is opened. [Both the tanks are open at top]\par
\qfig{p2_phy_s1_q01_1}
(viscosity of water $\eta$ = 1$\times$10$^{-3}$ Pa-s)\par
Given poiseuelle eqn for volume flow rate of fluid in a tube\par
$\dfrac{dV}{dt} = \dfrac{\Delta P\,\pi r^4}{8\eta L}$\par
$\Delta$P $\rightarrow$ pressure difference across ends of the tube\par
r $\rightarrow$ radius of tube\par
L $\rightarrow$ length of tube
\begin{qoptions}
\opt{A}{The difference between the levels of the water in the tanks decreases e times in time 1.6 $\times$ 10$^4$ s}
\opt{B}{The difference between the levels of the water in the tanks decreases e times in time 3.2 $\times$ 10$^4$ s}
\opt{C}{Total mechanical energy of the fluid is conserved during the process}
\opt{D}{Water level in the two containers cannot become equal in a finite time}
\end{qoptions}
\qn{2} A small charged bead can slide on a circular frictionless, insulating wire frame. A point like dipole is fixed at the centre of circle, dipole moment is $\vec{p}$. Initially the bead is on the plane of symmetry of the dipole. Bead is released from rest. Ignore the effect of gravity. Mark the correct options\par
\qfig{p2_phy_s1_q02_1}
\begin{qoptions}
\opt{A}{Magnitude of velocity of bead as function of its angular position is $\sqrt{\dfrac{\text{Qp cos}\,\theta}{2\pi\epsilon_0 \text{mr}^2}}$}
\opt{B}{Normal force exerted by the string on bead is zero at all points}
\opt{C}{If the wire frame were not present bead executes circular motion and returns to initial point after tracing a complete circle.}
\opt{D}{Bead would move along a circular path until it reached the opposite its starting position and then executes periodic motion}
\end{qoptions}
\qn{3} A point charge q = 6$\mu$C is moving in a straight line with a velocity $\vec{v} = 5\times10^4\,\hat{i}$ (m / s). When the charge is at the location P (3m, 4m, 0) choose the correct statements about the electric \& magnetic fields produced by the charge at the origin.
\begin{qoptions}
\opt{A}{Magnitude of magnetic field is 9.6 $\times$ 10$^{-10}$ T}
\opt{B}{Magnetic field is in --z direction}
\opt{C}{Electric field is varying with time}
\opt{D}{Magnetic field is decreasing (with time) in magnitude}
\end{qoptions}
\qn{4} Two cars X and Y are moving with speed 15 m/s and 11 m/s respectively in opposite directions approaching each other from far. The driver in car X blows a horn which has components of frequencies ranging from 650 Hz to 800 Hz. The band width of frequencies is thus 150 Hz, speed of sound is 340 m/s. The correct statement(s) for observer in car Y is : -
\begin{qoptions}
\opt{A}{The bandwidth of frequencies is 150 Hz}
\opt{B}{The bandwidth of frequencies is 162 Hz}
\opt{C}{Speed of sound of horn is 351 m/s}
\opt{D}{Speed of sound of horn is 329 m/s}
\end{qoptions}
\qn{5} Two beads of mass m are positioned at the top of a frictionless hoop of mass M and radius R, which stands vertically on ground. The beads are now given tiny kicks, and they slide down the hoop, one to the the right and one to the left, as shown in figure. Then choose the correct options\par
\qfig{p2_phy_s1_q05_1}
\begin{qoptions}
\opt{A}{The contact force from the ground immediately after the tiny kicks on the hoop is (M + 2m)g}
\opt{B}{The contact force from the ground on the hoop is (M + 2m)g at the instant when radius vector makes an angle cos$^{-1}$(2/3) with the upward vertical.}
\opt{C}{The smallest value of $\dfrac{\text{m}}{\text{M}}$ for which the hoop will rise off the ground at some time during the motion is $\dfrac{3}{2}$.}
\opt{D}{The hoop will never rise off the ground irrespective of the values of m and M.}
\end{qoptions}
\qn{6} Figure shows a two branched parallel circuit with R$_A$ = 10$\Omega$, L = $\dfrac{\sqrt{3}}{10}$H, R$_B$ = 20$\Omega$ and C = $\dfrac{\sqrt{3}}{2}$mH. Current in L -- R$_A$ is I$_1$ and in C -- R$_B$ is I$_2$ and main current is I\par
\qfig{p2_phy_s1_q06_1}
\begin{qoptions}
\opt{A}{Phase difference between I$_1$ and I$_2$ is 90$^\circ$}
\opt{B}{At some instant current in L-R$_A$ is 10A. At the same instant current in C-R$_B$ branch will be $5\sqrt{3}$A}
\opt{C}{At some instant I$_1$ is $10\sqrt{2}$ A then at this instant I will be $10\sqrt{2}$A}
\opt{D}{Power dissipated in the circuit is 2121.3 W}
\end{qoptions}
\sect{SECTION-I (ii)}
\qn{7} The arrangement is kept in a vertical plane and all masses are released from rest with strings taut. In list-I different values of masses are given. Match with corresponding parameters in list-II. Symbols have their usual meanings. Pulleys and strings are ideal.\par
\qfig{p2_phy_s2_q07_1}
{\footnotesize\setlength{\tabcolsep}{2pt}\noindent\begin{tabular}{|L{.7cm}|L{.6cm}|L{.6cm}|L{.6cm}|L{.6cm}|L{.7cm}|L{2cm}|}
\hline
\multicolumn{5}{|c|}{List-I} & \multicolumn{2}{c|}{List-II} \\ \hline
 & m$_1$ & m$_2$ & m$_3$ & m$_4$ & & \\ \hline
(P) & 2m & m & 3m & 4m & (1) & $|a_1| = |a_3|$ \\ \hline
(Q) & m & 2m & 4m & 3m & (2) & T$_1$ = T$_2$ \\ \hline
(R) & m & 2m & 3m & 4m & (3) & T$_2$ $>$ T$_1$ \\ \hline
(S) & 4m & 3m & 2m & m & (4) & $|a_1| > |a_3|$ \\ \hline
 & & & & & (5) & T$_3$ = 0 \\ \hline
\end{tabular}}\par
\begin{qoptions}
\opt{A}{P $\rightarrow$ 1,3;Q $\rightarrow$ 3,4;R $\rightarrow$ 4,5;S $\rightarrow$ 1,5}
\opt{B}{P $\rightarrow$ 1,3;Q $\rightarrow$ 3,4,5;R $\rightarrow$ 3,4,5;S $\rightarrow$ 5}
\opt{C}{P $\rightarrow$ 3,4;Q $\rightarrow$ 3,4;R $\rightarrow$ 1,5;S $\rightarrow$ 2,3,4}
\opt{D}{P $\rightarrow$ 3,4;Q $\rightarrow$ 2,4,5;R $\rightarrow$ 1,3;S $\rightarrow$ 4}
\end{qoptions}
\qn{8} Two light rods of length 1 m each are hinged together as shown in figure. Rod AB makes an angle $\theta$ with vertical while rod BC makes an angle $\phi$ with horizontal. End C of rod BC remains in contact with horizontal. Rod AB is rotated with constant angular velocity $\omega$ =1 rad/s in clockwise direction. At the instant when $\theta$ = 30$^\circ$ and $\phi$ = 30$^\circ$ match the variables in list-I with values in list-II.\par
\qfig{p2_phy_s2_q08_1}
{\footnotesize\setlength{\tabcolsep}{2pt}\noindent\begin{tabular}{|L{.6cm}|L{3cm}|L{.6cm}|L{2cm}|}
\hline
\multicolumn{2}{|c|}{List-I} & \multicolumn{2}{c|}{List-II} \\ \hline
(P) & Angular velocity of rod BC in rad/s & (1) & $\dfrac{(3\sqrt{3}+1)}{3\sqrt{3}}$ \\ \hline
(Q) & Velocity of block D in m/s & (2) & $\dfrac{(\sqrt{3}-1)}{\sqrt{6}}$ \\ \hline
(R) & Magnitude of angular acceleration of rod BC in rad/s$^2$ & (3) & $\dfrac{1}{\sqrt{3}}$ \\ \hline
(S) & Acceleration of point B in m/s$^2$ & (4) & 1 \\ \hline
\end{tabular}}\par
\begin{qoptions}
\opt{A}{P $\rightarrow$ 4;Q $\rightarrow$ 1;R $\rightarrow$ 2;S $\rightarrow$ 3}
\opt{B}{P $\rightarrow$ 1;Q $\rightarrow$ 4;R $\rightarrow$ 3;S $\rightarrow$ 2}
\opt{C}{P $\rightarrow$ 3;Q $\rightarrow$ 3;R $\rightarrow$ 1;S $\rightarrow$ 4}
\opt{D}{P $\rightarrow$ 3;Q $\rightarrow$ 2;R $\rightarrow$ 4;S $\rightarrow$ 1}
\end{qoptions}
\qn{9} Consider a system of five large conducting plates of area A. The charge on plate 1, 2, 3, 4, 5 are given as Q, 2Q, 3Q, 4Q and 5Q respectively. Initially the two switches are opened. Area of each plate is A. The distance between every successive plate is very very small equal to d.\par
\qfig{p2_phy_s2_q09_1}
{\footnotesize\setlength{\tabcolsep}{2pt}\noindent\begin{tabular}{|L{.6cm}|L{3.3cm}|L{.6cm}|L{2cm}|}
\hline
\multicolumn{2}{|c|}{List-I} & \multicolumn{2}{c|}{List-II} \\ \hline
(P) & The charge on surface b is & (1) & $-\dfrac{9}{2}\left(\dfrac{\text{Qd}}{\text{A}\epsilon_0}\right)$ \\ \hline
(Q) & The potential difference between plate 2 and 3 (V$_2$ -- V$_3$) is & (2) & --0.6 Q \\ \hline
\multicolumn{4}{|L{6.5cm}|}{Now both the switches S$_1$ and S$_2$ are closed simultaneously. Now Consider a new steady state.} \\ \hline
(R) & The charge on surface b is & (3) & --1.3 Q \\ \hline
(S) & The charge on surface f is & (4) & None of these \\ \hline
\end{tabular}}\par
\begin{qoptions}
\opt{A}{P $\rightarrow$ 4;Q $\rightarrow$ 3;R $\rightarrow$ 1;S $\rightarrow$ 4}
\opt{B}{P $\rightarrow$ 2;Q $\rightarrow$ 4;R $\rightarrow$ 3;S $\rightarrow$ 2}
\opt{C}{P $\rightarrow$ 3;Q $\rightarrow$ 3;R $\rightarrow$ 1;S $\rightarrow$ 4}
\opt{D}{P $\rightarrow$ 4;Q $\rightarrow$ 1;R $\rightarrow$ 3;S $\rightarrow$ 4}
\end{qoptions}
\qn{10} Phase space diagrams are useful tools in analyzing all kinds of dynamical problems. They are especially useful in studying the changes in motion as initial position and momentum are changed. Here we consider some simple dynamical systems for which phase space is a plane in which position is plotted along horizontal axis and momentum is plotted along vertical axis. The phase space diagram is x(t) vs. p(t) curve in this plane. The arrow on the curve indicates the time flow. For example, the phase space diagram for a particle moving with constant velocity is a straight line as shown in the figure. Similarly we may also plot momentum of a pendulum versus $\theta$ (with sign convention shown in figure (b)).\par
\qfig{p2_phy_s2_q10_1}
Figure (b) shows phase diagram of motion of simple pendulum (momentum P versus angle $\theta$). Choose potential energy level at the lowest point of the pendulum. E represents total energy of simple pendulum. Pendulum has a point mass connected with light rod.\par
\qfig{p2_phy_s2_q10_2}
{\footnotesize\setlength{\tabcolsep}{2pt}\noindent\begin{tabular}{|L{.6cm}|L{2cm}|L{.6cm}|L{3cm}|}
\hline
\multicolumn{2}{|c|}{List-I} & \multicolumn{2}{c|}{List-II} \\ \hline
(P) & Phase diagram a & (1) & E $<$ 2 mg$\ell$ \\ \hline
(Q) & Phase diagram b & (2) & E $\geq$ 2mg$\ell$ \\ \hline
(R) & Phase diagram c & (3) & May perform periodic and oscillatory motion \\ \hline
(S) & Phase diagram d & (4) & May represent SHM \\ \hline
 & & (5) & May represent angular velocity $\omega$ versus $\theta$ for a pendulum bob \\ \hline
\end{tabular}}\par
\begin{qoptions}
\opt{A}{P $\rightarrow$ 1,3,4,5;Q $\rightarrow$ 1,3,4,5;R $\rightarrow$ 2,5;S $\rightarrow$ 2,5}
\opt{B}{P $\rightarrow$ 1,3,4,5;Q $\rightarrow$ 1,4,5;R $\rightarrow$ 1,5;S $\rightarrow$ 1,5}
\opt{C}{P $\rightarrow$ 4,5;Q $\rightarrow$ 1,3,4,5;R $\rightarrow$ 1,4,5;S $\rightarrow$ 2,5}
\opt{D}{P $\rightarrow$ 2,4,5;Q $\rightarrow$ 3,4,5;R $\rightarrow$ 1,4;S $\rightarrow$ 1,3,5}
\end{qoptions}
\sect{SECTION-II}
\qn{1} In YDSE experiment two thin transparent sheets are used in front of the slits S$_1$ and S$_2$, one of the thin sheet has refractive index $\mu_1$ = 1.6 and other has $\mu_2$ = 1.4. If both sheets have average thickness $\left(\dfrac{t_1+t_2}{2}\right)$ the central maxima is at 5 mm from center O. Now the sheets are replaced by two sheets of same material of refractive index $\left(\dfrac{\mu_1+\mu_2}{2}\right)$ but having thickness t$_1$ and t$_2$ now central maxima is observed at a distance of 8mm from center O. For t$_1$ $>$ t$_2$ find the ratio $\dfrac{t_1}{t_2}$. Given d = 1mm, D = 1m
\qn{2} A tank with small orifice contains oil on top of water. It is immersed in a large tank of the same oil. Water flows through the orifice, determine the time at which flow stops\par
Density of oil = 800kg/m$^3$\par
Density of water = 1000kg/m$^3$\par
Initially the level of oil in both the tanks was same.\par
If t be the time in sec at which the flow stops then find the value of t$^2$\par
\qfig{p2_phy_s3_q02_1}
\qn{3} A highly conducting cylinder that has cross-sectional area of 100 cm$^2$ and 50 cm deep is filled with air at 21$^\circ$C and 1.00 atm figure (a). A 20kg piston is now lowered into the cylinder, compressing the air trapped inside figure (b). Finally a 80 kg man stands on the piston, further compressing the air, which remains at 21$^\circ$C figure (c). What temperature T (in $^\circ$C) should the gas be maintained to raise the piston and the man back to h$_i$. Take g = 10m/s$^2$\par
\qfig{p2_phy_s3_q03_1}
\qn{4} An object of mass m = $\sqrt{8}$ kg rests on an inclined plane that makes angle $\theta$ = 45$^\circ$ with the horizontal floor. What minimum force (in N), parallel to the base of the incline must be applied to the object in order to begin to move it along the plane parallel to the floor as shown? The coefficient of static friction between the object and the plane is $\mu_s$ = 1.25.\par
\qfig{p2_phy_s3_q04_1}
\qn{5} A 3.5V mobile phone battery can produce 1A of current for 1 hour. This can be charged using a square solar panel 25 cm on each side. Assuming an efficiency of 10\% and an incident solar power of 1kWm$^{-2}$ what time (in minutes) is needed to charge the battery?
\qn{6} A parallel plate condenser, with plate area A and distance between plates d, is filled with a medium whose permittivity varies as ;\par
$\epsilon(x) = \epsilon_0 + \beta\epsilon_0\,x/d$ \quad 0 $<$ x $<$ $\dfrac{d}{2}$\par
$\epsilon(x) = \epsilon_0 + \beta\epsilon_0\dfrac{(d-x)}{d}$ \quad $\dfrac{d}{2}$ $<$ x $<$ d.\par
x is the distance from one of the plates. For what value of $\beta$ would the capacity of the condenser be $\dfrac{1}{2\ell n2}$ times that when it is completely filled with a uniform dielectric having dielectric constant $\beta$ ?
\qn{7} Section of a refracting surface is shown in figure. The section is symmetrical about x axis. Refractive index of medium on right of the surface is n with respect to the medium on left of the surface. Parallel monochromatic light rays are incident on the surface as shown in figure and the refracted rays are focused at F at a distance f from origin O of the co-ordinate system shown in figure. The equation y(x) of shape of the section shown is found to be (n$^2$ -- 1) x$^2$ + n$^2$y$^2$ -- Bn (n--1) fx = 0. Find B.\par
\qfig{p2_phy_s3_q07_1}
\qn{8} In the determination of Young's modulus by using Searle's method, a wire of length L = 2m and diameter d = 0.5 mm is used. For a load M = 2.5 kg, an extension $\ell$ = 0.25 mm in the length of the wire is observed. Quantities d and $\ell$ are measured using a screw gauge and a micrometer, respectively. They have the same pitch of 0.5 mm. The number of divisions on their circular scale is 100. The length of wire, acceleration due to gravity and mass M are known exactly. The percentage error in measurement of the Young's modulus is\underline{\hspace{2cm}}
\partsec{2}{CHEMISTRY}
\sect{SECTION-I (i)}
\qn{1} Correct statement :
\begin{qoptions}
\opt{A}{Starch is the plant storage of carbohydrate.}
\opt{B}{Natural rubber is the trans-polymer.}
\opt{C}{Nylon show positive Lassigne test with Na metal.}
\opt{D}{Bakelite is phenol-acetaldehyde resin.}
\end{qoptions}
\qn{2} Consider the following values of I.E.(eV) for elements W and X:\par
{\footnotesize\setlength{\tabcolsep}{3pt}\noindent\begin{tabular}{llll l}
Element & I.E.$_1$ & I.E.$_2$ & I.E.$_3$ & I.E.$_4$ \\
W & 10.5 & 15.5 & 24.9 & 79.8 \\
X & 8 & 14.8 & 78.9 & 105.8 \\
\end{tabular}}\par
Other two element Y and Z have outer electronic configuration ns$^2$np$^4$ and ns$^2$np$^5$ respectively. Then according to given information which of the following compound(s) is/are not possible.
\begin{qoptions}
\opt{A}{W$_2$Y$_3$}
\opt{B}{X$_2$Y$_3$}
\opt{C}{WZ$_2$}
\opt{D}{XZ$_2$}
\end{qoptions}
\qn{3} (I) Na + O$_2$ (excess) $\xrightarrow{\Delta}$ major product\par
(II) K + O$_2$ (excess) $\xrightarrow{\Delta}$ ?\par
(III) O$_2$ + PtF$_6$ $\xrightarrow{\Delta}$ red compound prepared by Neil bartlett.\par
select correct for above reactions
\begin{qoptions}
\opt{A}{Electron population in $\pi^*$ molecular orbitals of O$_2$ is increases in (I) and (II)}
\opt{B}{Electron population in $\pi^*$ molecular orbitals of O$_2$ is decreases in (III)}
\opt{C}{$\pi$ bond order of O--O bond is lowest in product of IIIrd reaction}
\opt{D}{Fractinal O--O bond order is observed in product of Ist reaction.}
\end{qoptions}
\qn{4} Which of the following polymer is/are Biodegradable?
\begin{qoptions}
\opt{A}{\qfigi{p2_chem_s1_q04_A}}
\opt{B}{\qfigi{p2_chem_s1_q04_B}}
\opt{C}{\qfigi{p2_chem_s1_q04_C}}
\opt{D}{\qfigi{p2_chem_s1_q04_D}}
\end{qoptions}
\qn{5} Select correct statement(s)
\begin{qoptions}
\opt{A}{Plaster of paris is a hemihydrate of CaSO$_4$ obtained by heating the gypsum around 300 K.}
\opt{B}{Sodium carbonate is used in hard-water softening.}
\opt{C}{Sodium peroxide reacts with cold water giving hydrogen peroxide}
\opt{D}{The mobilities of the alkali metal ions in aqueous solutions are Li$^+$ $>$ Na$^+$ $>$ K$^+$ $>$ Rb$^+$ $>$ Cs$^+$}
\end{qoptions}
\qn{6} In a hypothetical solid, ``C'' atoms forming cubical closed packed lattice. ``A'' atoms occupy all tetrahedral void and ``B'' atoms occupy all octahedral voids. There is no distortion in ccp lattice. Fraction of body diagonal not covered up by atoms is : :
\begin{qoptions}
\opt{A}{0.76}
\opt{B}{0.24}
\opt{C}{0.68}
\opt{D}{0.12}
\end{qoptions}
\sect{SECTION-I (ii)}
\qn{7} {\footnotesize\setlength{\tabcolsep}{2pt}\noindent\begin{tabular}{|L{.6cm}|L{2.2cm}|L{.6cm}|L{2.6cm}|}
\hline
\multicolumn{2}{|c|}{List-I} & \multicolumn{2}{c|}{List-II} \\ \hline
(P) & \qfigi{p2_chem_s2_q07_1} & (1) & Disproportionation reaction with OH$^-$ \\ \hline
(Q) & \qfigi{p2_chem_s2_q07_2} & (2) & Condensation reaction with OH$^-$ \\ \hline
(R) & \qfigi{p2_chem_s2_q07_3} & (3) & Silver mirror test \\ \hline
(S) & \qfigi{p2_chem_s2_q07_4} & (4) & Haloform test \\ \hline
\end{tabular}}\par
\begin{qoptions}
\opt{A}{P $\rightarrow$ 1,3;Q $\rightarrow$ 2,4;R $\rightarrow$ 1,3;S $\rightarrow$ 2,3}
\opt{B}{P $\rightarrow$ 2,3,4;Q $\rightarrow$ 1,3;R $\rightarrow$ 2,4;S $\rightarrow$ 1,3}
\opt{C}{P $\rightarrow$ 3,4;Q $\rightarrow$ 2,4;R $\rightarrow$ 1,2,3;S $\rightarrow$ 1,2,3}
\opt{D}{P $\rightarrow$ 1,2,3;Q $\rightarrow$ 2,3;R $\rightarrow$ 3,4;S $\rightarrow$ 2,3}
\end{qoptions}
\qn{8} {\footnotesize\setlength{\tabcolsep}{2pt}\noindent\begin{tabular}{|L{.6cm}|L{1cm}|L{.6cm}|L{3.8cm}|}
\hline
\multicolumn{2}{|c|}{List-I} & \multicolumn{2}{c|}{List-II} \\ \hline
(P) & Fe$^{n+}$ & (1) & (For n = 0), forms white ppt with dil. HCl. \\ \hline
(Q) & Cu$^{n+}$ & (2) & (For n = 2), forms coloured ppt with excess of NaOH solution. \\ \hline
(R) & Hg$^{n+}$ & (3) & (For n = 0), N$_2$O gas will be released with dil.HNO$_3$. \\ \hline
(S) & Pb$^{n+}$ & (4) & (For n = 2), Forms soluble complex with excess of NH$_3$ solution. \\ \hline
 & & (5) & (For n = 2) forms black ppt on passing H$_2$S gas in its aqueous solution. \\ \hline
\end{tabular}}\par
\begin{qoptions}
\opt{A}{P $\rightarrow$ 2,3;Q $\rightarrow$ 2,4,5;R $\rightarrow$ 2,5;S $\rightarrow$ 1,5}
\opt{B}{P $\rightarrow$ 1,2,3;Q $\rightarrow$ 3,5;R $\rightarrow$ 1,2,5;S $\rightarrow$ 2,4,5}
\opt{C}{P $\rightarrow$ 1,4;Q $\rightarrow$ 2,4;R $\rightarrow$ 3,5;S $\rightarrow$ 1,3}
\opt{D}{P $\rightarrow$ 2,4,5;Q $\rightarrow$ 2,3,5;R $\rightarrow$ 1,5;S $\rightarrow$ 1,4,5}
\end{qoptions}
\qn{9} Complete the following reactions identify the major products and apply a chemical test to distinguish between following pairs :\par
{\footnotesize\setlength{\tabcolsep}{2pt}\noindent\begin{tabular}{|L{.6cm}|L{2cm}|L{.6cm}|L{2.4cm}|}
\hline
\multicolumn{2}{|c|}{List-I (Pairs)} & \multicolumn{2}{c|}{List-II (Chemical Test)} \\ \hline
(P) & J and M & (1) & Na-metal \\ \hline
(Q) & J and K & (2) & NaOI \\ \hline
(R) & K and M & (3) & 2,4-DNP \\ \hline
(S) & L and M & (4) & NaHSO$_3$ \\ \hline
\end{tabular}}\par
Reactions are as follows\par
\qfig{p2_chem_s2_q09_1}
\begin{qoptions}
\opt{A}{P $\rightarrow$ 3;Q $\rightarrow$ 2;R $\rightarrow$ 1;S $\rightarrow$ 4}
\opt{B}{P $\rightarrow$ 2;Q $\rightarrow$ 3;R $\rightarrow$ 4;S $\rightarrow$ 1}
\opt{C}{P $\rightarrow$ 2;Q $\rightarrow$ 1;R $\rightarrow$ 3;S $\rightarrow$ 4}
\opt{D}{P $\rightarrow$ 1;Q $\rightarrow$ 2;R $\rightarrow$ 4;S $\rightarrow$ 3}
\end{qoptions}
\qn{10} {\footnotesize\setlength{\tabcolsep}{2pt}\noindent\begin{tabular}{|L{.6cm}|L{3.6cm}|L{.6cm}|L{2cm}|}
\hline
\multicolumn{2}{|c|}{List-I} & \multicolumn{2}{c|}{List-II} \\ \hline
(P) & \qfigi{p2_chem_s2_q10_1} & (1) & Carbanion is intermediate \\ \hline
(Q) & \qfigi{p2_chem_s2_q10_2} & (2) & Carbene is intermediate \\ \hline
(R) & \qfigi{p2_chem_s2_q10_3} & (3) & Product containing N- atom \\ \hline
(S) & \qfigi{p2_chem_s2_q10_4} & (4) & Cyanide or isocyanide is product \\ \hline
\end{tabular}}\par
\begin{qoptions}
\opt{A}{P $\rightarrow$ 1,3,4;Q $\rightarrow$ 3,4;R $\rightarrow$ 2,3;S $\rightarrow$ 2,3,4}
\opt{B}{P $\rightarrow$ 1,2,3,4;Q $\rightarrow$ 3;R $\rightarrow$ 1,2;S $\rightarrow$ 3,4}
\opt{C}{P $\rightarrow$ 2;Q $\rightarrow$ 1,3,4;R $\rightarrow$ 2,3;S $\rightarrow$ 1,2}
\opt{D}{P $\rightarrow$ 3,4;Q $\rightarrow$ 3,4;R $\rightarrow$ 2,3;S $\rightarrow$ 2,3}
\end{qoptions}
\sect{SECTION-II}
\qn{1} Find total number of species in which incoming electron enters into H.O.M.O.\par
H$_2$, Li$_2$, B$_2$, C$_2$, N$_2$, O$_2$, F$_2$, Na$_2$, S$_2$, Cl$_2$.
\qn{2} 40 ml of 0.05 M solution of sodium sesquicarbonate (Na$_2$CO$_3$.NaHCO$_3$. 2H$_2$O) is titrated against 0.05 M HCl. ``X'' ml of HCl is used when phenolphthalein is the indicator and ``Y'' ml of HCl is used when methyl orange is the indicator in two separate titrations. Hence, value of $\dfrac{(Y-X)}{10}$ is
\qn{3} Following reaction takes place at $\left(\dfrac{600}{2.303\times0.8314}\right)$ K temperature.\par
H$_2$(g) + 2Ag$^+$(aq) $\rightleftharpoons$ 2Ag(s) + 2H$^+$(aq)\par
P$_{\text{H}_2}$ = 1.0 bar, [Ag$^+$] = 10$^{-6}$ M, [H$^+$] = 10$^{-3}$ M , \quad $\Delta_f$G$^\circ$(Ag$^+$,aq) = 75 kJ mol$^{-1}$\par
Calculate z. where z = $\dfrac{|\Delta_r\text{G(in kJ)}|}{19}$
\qn{4} Select the number of ores for which roasting process is applied during metallurgy.\par
Argentite, Copper pyrites, Galena, Zinc blende, Chalcocite
\qn{5} In the saturated aqueous solution of PbCl$_2$ the freezing point decreases by $\left(\dfrac{X}{100}\right)^\circ$C then `X' is :\par
(Given k$_{sp}$ of PbCl$_2$ = 4 $\times$ 10$^{-6}$, K$_{f,\text{water}}$ = 2 k kg/mole)
\qn{6} Total number of compounds which are soluble in a hot aqueous NaOH are:\par
(i) Salicyclic acid\par
(ii) Aspirine\par
(iii) Formic acid\par
(iv) Acetic acid\par
(v) Succinic anhydride\par
(vi) Cyclohexanone\par
(vii) Benzene sulphonamide\par
(viii) Cyclohexene
\qn{7} In how many compounds intermolecular H--bonding will exist.\par
\qfig{p2_chem_s3_q07_1}
\qn{8} A sample of H-atoms containing all the atoms in a particular excited state, absorb radiations of a particular wave length by which the atoms get excited to another excited state. When the atoms finally de-excite to the ground state, they emit radiations of 15 different wavelength. Out of these 15 radiations, 5 have wavelengths shorter than the absorbed radiation and 9 have wavelength longer than the absorbed radiation. What is the initial excited state of atoms?
\partsec{3}{MATHEMATICS}
\sect{SECTION-I (i)}
\qn{1} Let f(x) = 2x$^3$ -- 15x$^2$ + 36x -- 23 and g(x) = $\begin{cases}\max\{f(t),\ 1\leqslant t\leqslant x\}, & \text{if } 1\leqslant x\leqslant\frac{7}{2}\\ 12-2x, & \text{if } \frac{7}{2}<x\leqslant 6\end{cases}$\par
Which is correct about g(x)
\begin{qoptions}
\opt{A}{Continuous everywhere}
\opt{B}{Differentiable everywhere}
\opt{C}{has greatest value 5}
\opt{D}{None of these}
\end{qoptions}
\qn{2} For distinct complex numbers z$_1$, z$_2$, ..., z$_n$, (n $>$ 2) the value of $\dfrac{|z_2-z_1|^2+|z_3-z_2|^2+......+|z_n-z_{n-1}|^2}{|z_n-z_1|^2}$ can not be less than
\begin{qoptions}
\opt{A}{$\dfrac{1}{n-1}$}
\opt{B}{$\dfrac{1}{n}$}
\opt{C}{$\dfrac{1}{n+1}$}
\opt{D}{$\dfrac{1}{n+2}$}
\end{qoptions}
\qn{3} Equation(s) of possible common tangents to y$^2$ = 8(x -- 3) and x$^2$ = 8(y -- 3) is/are :
\begin{qoptions}
\opt{A}{x + y = 1}
\opt{B}{2x -- y = 5}
\opt{C}{x -- 2y + 5 = 0}
\opt{D}{None of these}
\end{qoptions}
\qn{4} If the line $\dfrac{x}{1}=\dfrac{y}{2}=\dfrac{z}{3}$ intersects line 3$\beta^2$x + 3(1 -- 2$\alpha$)y + z = 3 = $-\dfrac{1}{2}$(6$\alpha^2$x + 3(1 -- 2$\beta$)y + 2z) then point ($\alpha$, $\beta$, 1) lies on the plane
\begin{qoptions}
\opt{A}{2x -- y + z = 4}
\opt{B}{x + y -- z = 2}
\opt{C}{x -- 2y = 0}
\opt{D}{2x -- y = 0}
\end{qoptions}
\qn{5} Let $\displaystyle\int_0^\infty e^{-x^2}dx = a$ and I$_n$ = $\displaystyle\int_0^\infty x^n e^{-x^2}dx$, n $\in$ N, then
\begin{qoptions}
\opt{A}{I$_n$ = $\left(\dfrac{n-1}{2}\right)$ I$_{n-2}$}
\opt{B}{I$_n$ = $\dfrac{(n-1)!}{2^{n-1}\left(\frac{n-2}{2}\right)!}$ a, if n is even}
\opt{C}{I$_n$ = $\dfrac{1}{2}\left(\dfrac{n-2}{2}\right)!$ a, if n is even}
\opt{D}{I$_n$ = $\dfrac{(n-1)!}{2^{n-1}\left(\frac{n-2}{2}\right)!}$, if n is even}
\end{qoptions}
\qn{6} In $\Delta$ABC (with usual notation), if cos A + cos B = 4sin$^2\dfrac{C}{2}$, then which of the following hold(s) good ?
\begin{qoptions}
\opt{A}{cot$\dfrac{A}{2}$cot$\dfrac{B}{2}$ = 2}
\opt{B}{cot$\dfrac{A}{2}$cot$\dfrac{B}{2}$ = 3}
\opt{C}{a, c, b are in A.P.}
\opt{D}{a, b, c are in G.P.}
\end{qoptions}
\sect{SECTION-I (ii)}
\qn{7} Match the following Column-I with Column-II\par
{\footnotesize\setlength{\tabcolsep}{2pt}\noindent\begin{tabular}{|L{.6cm}|L{3.8cm}|L{.6cm}|L{2cm}|}
\hline
\multicolumn{2}{|c|}{List-I} & \multicolumn{2}{c|}{List-II} \\ \hline
(P) & If A = [a$_{ij}$]$_{3\times3}$ and a$_{ij}$ = i$^2$ + j$^2$, then A is & (1) & singular \\ \hline
(Q) & A = [a$_{ij}$]$_{3\times3}$ and a$_{ij}$ = 3$^{i-j}$, then A is & (2) & skew-symmetric \\ \hline
(R) & A = [a$_{ij}$]$_{3\times3}$ and a$_{ij}$ = i$^2$ -- j$^2$, then A is & (3) & symmetric \\ \hline
(S) & A = $\begin{bmatrix}2&-2&-4\\-1&3&4\\1&-2&-3\end{bmatrix}$, then A is & (4) & idempotent \\ \hline
\end{tabular}}\par
\begin{qoptions}
\opt{A}{P $\rightarrow$ 3;Q $\rightarrow$ 1;R $\rightarrow$ 1;S $\rightarrow$ 2}
\opt{B}{P $\rightarrow$ 3;Q $\rightarrow$ 2;R $\rightarrow$ 1;S $\rightarrow$ 1}
\opt{C}{P $\rightarrow$ 1;Q $\rightarrow$ 2;R $\rightarrow$ 3;S $\rightarrow$ 4}
\opt{D}{P $\rightarrow$ 3;Q $\rightarrow$ 1;R $\rightarrow$ 2;S $\rightarrow$ 1}
\end{qoptions}
\qn{8} Column-I with Column-II (where [.] denotes GIF and \{.\} denotes fractional part function)\par
{\footnotesize\setlength{\tabcolsep}{2pt}\noindent\begin{tabular}{|L{.6cm}|L{4.6cm}|L{.6cm}|L{.6cm}|}
\hline
\multicolumn{2}{|c|}{List-I} & \multicolumn{2}{c|}{List-II} \\ \hline
(P) & For f(x) = $\begin{cases}\dfrac{|x-2|}{x-2}[\log_e x], & 1\leqslant x\leqslant 3-\{2\}\\ \{x^2\}, & 3<x\leqslant\frac{7}{2}\end{cases}$ number of points of discontinuity is & (1) & 0 \\ \hline
(Q) & For f(x) = $\begin{cases}\left(\sqrt{2}+\sin\frac{1}{x}\right)e^{\frac{-1}{|x|}}, & \text{if } x\neq0\\ 0, & \text{if } x=0\end{cases}$ number of points of extremum is & (2) & 5 \\ \hline
(R) & number of points where f(x) = max $\{|3-x^2|, 3-x^3\}$ is not differentiable & (3) & 2 \\ \hline
(S) & $2\displaystyle\lim_{x\to0}x^2\left(1+2+3+....+\left[\dfrac{1}{|x|}\right]\right)$ =, where [x] represents GIF & (4) & 1 \\ \hline
\end{tabular}}\par
\begin{qoptions}
\opt{A}{P $\rightarrow$ 2;Q $\rightarrow$ 1;R $\rightarrow$ 4;S $\rightarrow$ 3}
\opt{B}{P $\rightarrow$ 2;Q $\rightarrow$ 3;R $\rightarrow$ 1;S $\rightarrow$ 4}
\opt{C}{P $\rightarrow$ 2;Q $\rightarrow$ 4;R $\rightarrow$ 3;S $\rightarrow$ 4}
\opt{D}{P $\rightarrow$ 1;Q $\rightarrow$ 2;R $\rightarrow$ 3;S $\rightarrow$ 4}
\end{qoptions}
\qn{9} Match the following :\par
{\footnotesize\setlength{\tabcolsep}{2pt}\noindent\begin{tabular}{|L{.6cm}|L{3.6cm}|L{.6cm}|L{1.6cm}|}
\hline
\multicolumn{2}{|c|}{List-I} & \multicolumn{2}{c|}{List-II} \\ \hline
(P) & If $\sin^{-1}|\cos x| - \cos^{-1}|\sin x|$ = a has atleast one solution, then the exhaustive set in which `a' can lie & (1) & \{0\} \\ \hline
(Q) & If equation $\sin^{-1}\sqrt{x}+\cos^{-1}\sqrt{x^2-1}+\tan^{-1}\tan y$ = a has atleast one solution, then the exhaustive set in which `a' can lie & (2) & $\left(\dfrac{\pi}{2},\dfrac{3\pi}{2}\right)$ \\ \hline
(R) & If $\sqrt{|\sin^{-1}|\sin x||}+\sqrt{|\tan^{-1}|\tan x||}$ = $2\sqrt{a}$ has atleast one solution, then the exhaustive set in which `a' can lie & (3) & $\left[0,\dfrac{\pi}{2}\right)$ \\ \hline
(S) & $\sin^{-1}(x^2+y^2)+\tan^{-1}\sqrt{4y^2-1}+\sec^{-1}x$ = a, then the exhaustive set in which `a' can lie & (4) & $\phi$ \\ \hline
\end{tabular}}\par
\begin{qoptions}
\opt{A}{P $\rightarrow$ 1;Q $\rightarrow$ 2;R $\rightarrow$ 4;S $\rightarrow$ 3}
\opt{B}{P $\rightarrow$ 1;Q $\rightarrow$ 4;R $\rightarrow$ 3;S $\rightarrow$ 2}
\opt{C}{P $\rightarrow$ 2;Q $\rightarrow$ 1;R $\rightarrow$ 4;S $\rightarrow$ 3}
\opt{D}{P $\rightarrow$ 1;Q $\rightarrow$ 2;R $\rightarrow$ 3;S $\rightarrow$ 4}
\end{qoptions}
\qn{10} Column-I with Column-II\par
{\footnotesize\setlength{\tabcolsep}{2pt}\noindent\begin{tabular}{|L{.6cm}|L{3.4cm}|L{.6cm}|L{2.2cm}|}
\hline
\multicolumn{2}{|c|}{Column-I} & \multicolumn{2}{c|}{Column-II} \\ \hline
(P) & Vertices of a triangle are (1, 2), (3, 4) \& (2, --4) & (1) & $\lambda\neq-7$ \\ \hline
(Q) & Equation of sides of a triangle are x + y = 3, 3x -- 5y + 7 = 0 and x -- 3y + 1 = 0 & (2) & triangle is obtuse angled \\ \hline
(R) & Orthocenter and circumcenter are on different half plane w.r.t larger side. & (3) & (1, 1) lies inside the triangle \\ \hline
(S) & Equations of sides of a triangle are, 3x -- 4y -- 13 = 0, 8x -- 11y -- 33 = 0 \& 2x -- 3y + $\lambda$ = 0 & (4) & centroid lies on 3y = 2 \\ \hline
\end{tabular}}\par
\begin{qoptions}
\opt{A}{P $\rightarrow$ 2;Q $\rightarrow$ 3;R $\rightarrow$ 2;S $\rightarrow$ 4}
\opt{B}{P $\rightarrow$ 2,4;Q $\rightarrow$ 2,3;R $\rightarrow$ 2;S $\rightarrow$ 1,4}
\opt{C}{P $\rightarrow$ 2,4;Q $\rightarrow$ 2,3,4;R $\rightarrow$ 2;S $\rightarrow$ 1,2}
\opt{D}{P $\rightarrow$ 2,4;Q $\rightarrow$ 4;R $\rightarrow$ 2;S $\rightarrow$ 1,2,4}
\end{qoptions}
\sect{SECTION-II}
\qn{1} The value of the expression ${}^{n+1}C_2$ + 2$\left({}^2C_2+{}^3C_2+{}^4C_2+.......+{}^nC_2\right)$, for n = 6 is
\qn{2} A sequence is obtained by deleting all perfect squares from set of natural number. The remainder when the 2003$^{\text{rd}}$ term of new sequence is divided by 2048, is ?
\qn{3} Let 2(f(x))$^2$ -- $\dfrac{d^2f(x)}{dx^2}$ f(x) + $\left(\dfrac{df(x)}{dx}\right)^2$ = 0 \& f(0) = f(1) = --1. Area of region bounded by y = 0, x = 0, x = 1 and y = (2x --1) f(x) is 2$\left(\dfrac{e^{1/a}-1}{e^{1/a}}\right)$, then `a' is ?
\qn{4} Two boys A and B find the jumble of n ropes lying on the floor. Each takes hold of one loose end randomly. If the probability that they are both holding the same rope is $\dfrac{1}{101}$ then the number of ropes is equal to ?
\qn{5} A circle of radius r = 4 units is inscribed in an equilateral triantgle ABC, then an equilateral triangle is inscribed in the circle, a circle again is inscribed in the later triangle and so on. In this way the process continues infinitely. If r, x$_1$, x$_2$, ..., x$_n$, .... be the radii of these circles respectively, then the sum of radii of all the circles x$_1$ + x$_2$ + x$_3$ ...$\infty$ is equals to
\qn{6} A trapezium is inscribed in the parabola y$^2$ = 4x such that its digonals pass through the focus and are of length $\dfrac{25}{4}$ units. If area of this trapezium is A, then 4A is ?
\qn{7} If $\alpha$, $\beta$ are the roots of the quadratic equation, x$^2$ -- 2p(x -- 4) -- 15 = 0, then the greatest integer p for which one root is less than 1 and the other root is greater than 2 is ?
\qn{8} Let $\vec{A}$, $\vec{B}$, $\vec{C}$ be vectors of length 3, 4, 5 respectively. Let $\vec{A}$ be perpendicular to $\vec{B}+\vec{C}$, $\vec{B}$ is perpendicular to $\vec{C}+\vec{A}$ and $\vec{C}$ is perpendicular to $\vec{A}+\vec{B}$ and if the length of vector $|\vec{A}+\vec{B}+\vec{C}|$ is k$\sqrt{2}$ then find the value of k
\end{multicols}
\end{document}
